\section{Порівняння криптографічних бібліотек}

\subsection{Обрані криптографічні бібліотеки}

% Коротко описати cryptography та PyCryptodome

\subsection{Порівняльний аналіз бібліотек}

\begin{longtable}{|p{2cm}|p{6cm}|p{6cm}|}
\hline
\textbf{Критерій} &
\textbf{1} &
\textbf{2} \\
\hline

Мови та платформи &
&
\\
\hline

Симетрич\-не шифрування &
&
\\
\hline

AEAD &
&
\\
\hline

Асимет\-рична криптографія &
&
\\
\hline

Цифровий підпис &
&
\\
\hline

Геш-функції та KDF &
&
\\
\hline

Ініціаліза\-ція генератора ви\-падкових чисел &
&
\\
\hline

Крипто\-графічно стійкі випадкові значення &
&
\\
\hline

Високо\-рівневий API &
&
\\
\hline

Апаратне прискорення &
&
\\
\hline

FIPS-сумісність / сертифікація &
&
\\
\hline

Ліцензія &
&
\\
\hline

Вразли\-вості та історія зламів &
&
\\
\hline

Підтримка спільнотою &
&
\\
\hline

Актуаль\-ність документації &
&
\\
\hline

\end{longtable}

\subsection{Обґрунтування вибору бібліотеки}

% Чому обираємо cryptography